\documentclass[10pt,a4paper]{amsart}
\usepackage[margin=19mm]{geometry}
\usepackage{amsmath,amssymb,mathtools}
\usepackage[scr=rsfs,cal=euler]{mathalfa}
\usepackage{xfrac}
\usepackage[unicode,hidelinks,bookmarks=false]{hyperref}
\newcommand{\ie}{i.e.\ }
\newcommand{\cf}{cf.\ }
\newcommand{\resp}{respectively\ }
\newcommand{\Subsection}{\subsection}
\providecommand{\nobreakdash}{\nobreak\hbox{-}}
\setlength{\emergencystretch}{2em}
\multlinegap0pt
\usepackage{amssymb}
\usepackage{mathtools}
\usepackage{xcolor}
\usepackage{dsfont}
\usepackage[shortlabels]{enumitem}
\newtheorem{lemma}{Lemma}
\newcommand{\gradgr}{\nabla_{\gamma}}
\newcommand{\om} {\Omega}
\newcommand{\spacegr}{H^{1}_{0,\gamma}(\Omega)}
\title{Existence and multiplicity of solutions for a critical Grushin problem with a singular nonlinearity}
\author{Shammi Malhotra}
\date{Source: arXiv:2605.02020v1 (CC BY 4.0)}
\begin{document}
\maketitle

\noindent\emph{Source and licence.} Existence and multiplicity of solutions for a critical Grushin problem with a singular nonlinearity. Version arXiv:2605.02020v1 (CC BY 4.0), distributed by arXiv under the Creative Commons Attribution 4.0 International licence, http://creativecommons.org/licenses/by/4.0/, as recorded in the arXiv licence metadata for this exact version and shown on the abstract page https://arxiv.org/abs/2605.02020v1. Full licence notice: \texttt{licenses/arxiv-2605.02020-CC-BY-4.0.txt}.\par\smallskip

\noindent\textit{Editorial scope.} Counted unit: the article's lemma lem:variational-identity (source0.tex lines 1783--1812). The article's complete original proof of it is delivered with it (source0.tex lines 1813--1949, one \texttt{proof} environment of 137 lines). No mathematics is written, repaired or re-derived: the statement and the proof are the article's own lines, in the article's own notation, and no retained line is substituted or re-flowed.  Retained uncounted as the source-local context the proof needs: lines 1738--1741.  Nothing was omitted from any retained span, so the package carries no omission record.  No retained line contains a citation, so the wrapper carries no bibliography and no bibliography entry.  No retained line contains a figure environment, an image-inclusion command or drawing code, so no image is delivered and the package declares no asset dependency.  Editorial formatting only, in addition: an AMS \texttt{amsart} preamble that restates the article's own declarations and macros used by the retained lines, namely the environments \texttt{lemma} on separate counters, together with the standard packages those lines need.  The retained environments print in this document as Lemma 1 in order of appearance, whereas the article prints the same items with the numbering of its own full document.  The complete native source of the article as distributed in the arXiv e-print for this exact version is bundled unchanged as \texttt{sources/199683/source0.tex}.
\begin{equation}
\label{eq:main_eq_singular}
- \Delta_{\gamma} u = g(z,u) + w \qquad \text{in } \Omega,
\end{equation}

\begin{lemma}
\label{lem:variational-identity}
Let $g : \Omega \times \mathbb{R} \to \mathbb{R}$ be a Carath\'{e}odory function such that for all $S > 0$,
\begin{equation*}
% \label{eq:g-growth}
\sup_{|s|\le S} |g(\cdot,s)| \in L^1_{loc}(\Omega).
\end{equation*}
Let $w \in H^{-1,2}_{\gamma}(\Omega)$, and $\varphi,\, u,\, \psi \in  H^{1}_{\gamma, loc}(\Omega)$. Assume that
\begin{itemize}
\item $\varphi$ is a subsolution of equation \eqref{eq:main_eq_singular};
\item $\psi$ is a supersolution of equation \eqref{eq:main_eq_singular};
\item$\varphi \le u \le \psi$ a.e. in $\om$;
\item$g(z,u) \in L^1_{loc}(\Omega)$;
\item for every $v \in u + \bigl( \spacegr \cap L^\infty_{loc}(\Omega) \bigr)$ with $\varphi \le v \le \psi$ a.e. in $\om$, the following holds
\begin{equation}\label{ineq:min_on_convex}
\int_\Omega
\gradgr u \cdot \gradgr (v-u)\,dz
\ge
\int_\Omega
g(z,u)(v-u)\,dz
+
\langle w , v-u \rangle.
\end{equation}
\end{itemize}
Then $u$ satisfies
\begin{equation*}
- \Delta_{\gamma} u = g(z,u) + w
\end{equation*}
in the sense of distributions. 
\end{lemma}

\begin{proof}
Let $\eta \in C_c^\infty(\mathbb{R})$ be such that
\begin{equation*}
0 \le \eta \le 1 \quad \text{in } \mathbb{R}, \qquad
\eta = 1 \ \text{on } [-1,1], \qquad
Supp(\eta) \subset (-2,2).
\end{equation*}
Let $v \in C_c^\infty(\Omega)$ with $v \ge 0$. Fix $k \ge 1$ and $t > 0$, and set
\begin{equation*}
v_k = \eta\!\left( \frac{u}{k} \right) v,
\qquad
v_{k,t} = \min \{ u + t v_k , \psi \}.
\end{equation*}
Since $u \le v_{k,t} \le \psi$, using $v_{k,t}$ as the test function in \eqref{ineq:min_on_convex}, we obtain
\begin{align*}
&\int_\Omega
\Big(
|\gradgr (v_{k,t}-u)|^2
-
\big( g(z,v_{k,t}) - g(z,u) \big)
(v_{k,t}-u)
\Big)\,dz
 \\ &
\le
\int_\Omega
\gradgr v_{k,t} \cdot \gradgr (v_{k,t}-u)
-
g(z,v_{k,t})(v_{k,t}-u)\,dz
-
\langle w , v_{k,t}-u \rangle
\\ & =
\int_\Omega
\gradgr v_{k,t} \cdot \gradgr (v_{k,t}-u-t v_k)
-
g(z,v_{k,t})(v_{k,t}-u-t v_k)\,dz
-
\langle w , v_{k,t}-u-t v_k \rangle
\\
&\quad
+ t \int_\Omega
\big(
\gradgr v_{k,t} \cdot \gradgr v_k
-
g(z,v_{k,t}) v_k
\big)\,dz
-
t \langle w , v_k \rangle \\
& =
\int_\Omega
\gradgr \psi \cdot \gradgr (v_{k,t}-u-t v_k)
-
g(z,\psi)(v_{k,t}-u-t v_k)\,dz
-
\langle w , v_{k,t}-u-t v_k \rangle
\\
&+ t \int_\Omega
\big(
\gradgr v_{k,t} \cdot \gradgr v_k
-
g(z,v_{k,t}) v_k
\big)\,dz
-
t \langle w , v_k \rangle 
\\& \leq 
t \int_\Omega
\big(
\gradgr v_{k,t} \cdot \gradgr v_k
-
g(z,v_{k,t}) v_k
\big)\,dz
-
t \langle w , v_k \rangle.
\end{align*}
where the last inequality follows from the fact that $\psi$ is a supersolution of equation \eqref{eq:main_eq_singular}.

Thus, the relation $v_{k,t} - u \le t v_k \le t v$ yields
\begin{equation*}
\int_\Omega
\big(
\gradgr v_{k,t} \cdot \gradgr v_k
-
g(z,v_{k,t}) v_k
\big)\,dz
-
\langle w , v_k \rangle
\ge
-
\int_\Omega
\big| g(z,v_{k,t}) - g(z,u) \big|\, |v_k|\, dz .
\end{equation*}

Since
\begin{equation*}
|g(z,v_{k,t})|\,|v_k|
\le
\left(
\sup_{|s|\le 2k + \|v\|_\infty}
|g(z,s)|
\right)
|v|,
\end{equation*}
we may pass to the limit as $t \to 0^{+}$ by the Dominated Convergence Theorem. Consequently, we obtain
\begin{equation*}
\int_\Omega
\big(
\gradgr u \cdot \gradgr v_k
-
g(z,u) v_k
\big)\,dz
-
\langle w , v_k \rangle
\ge 0 .
\end{equation*}
Letting $k \to \infty$, it follows that
\begin{equation}
\label{eq:ineq-ten}
\int_\Omega
\big(
\gradgr u \cdot \gradgr v
-
g(z,u) v
\big)\,dz
-
\langle w , v \rangle
\ge 0 ,
\end{equation}
for every $v \in C_c^\infty(\Omega)$ with $v \ge 0$. Now let $v \in C_c^\infty(\Omega)$ with $v \le 0$. For $k \ge 1$ and $t > 0$, define
\begin{equation*}
v_k = \eta\!\left( \frac{u}{k} \right) v,
\qquad
v_{k,t} = \max \{ u + t v_k , \varphi \}.
\end{equation*}
Repeating the previous argument, we deduce that inequality
\eqref{eq:ineq-ten} holds for every $v \in C_c^\infty(\Omega)$.

Finally, replacing $v$ by $-v$ in \eqref{eq:ineq-ten}, we conclude the proof.
\end{proof}

\end{document}
